\documentclass[parskip=full]{scrartcl}
\usepackage[top=2.5cm, bottom=2.5cm, left=2.5cm, right=2.5cm]{geometry}
\usepackage[utf8]{inputenc}
\usepackage[T1]{fontenc}
\usepackage[german]{babel}
\usepackage{hyperref}
\usepackage[toc, nonumberlist]{glossaries}
\usepackage{graphicx}
\usepackage{enumitem}
\usepackage{float}
\usepackage{color}

\title{OSIP - Softwaredesign}
\subtitle{OPC UA Simulator for Industrial Plants}
\author{
    M. Armbruster\\
    D. Kahles\\
    H. Lehmann\\
    M. Schwarzmann\\
    N. Wilhelm
}

\begin{document}
\maketitle

\vspace{20px}
\begin{center}
  \includegraphics[scale=0.4]{../icon.png}
\end{center}
\pagebreak
\tableofcontents
\pagebreak

\section{Einleitung}
% Aufteilung in die Pakete erklären

\section{Klassenbeschreibungen}
% JavaDoc

\section{Abläufe}
% Sequenzdiagramme

\section{Änderungen zum Pflichtenheft}
%Implementiert:
%- Logging von Fehler bei der Ausführung stdout
%- Szenarien
%- Jitter
%- Zeitintervall
%- Alarme ein-/ausschalten (und das abspeichern)
%- Logging für Alarme in der Überwachungskonsole
%
%nicht implementiert:
%- andere Sprache
%- Alarmschwellen konfigurieren
%- Jitter einstellen

\section{Klassendiagramme}
Im Anhang.

\section{Verwendete Entwurfsmuster}
% Aufzählen: Geheimnisprinzip, MVC etc.

\end{document}
